\scanpage{70}
% archive sheet number: 69
% L1-2
Supposons que $k'/k$ \unsure{soit|est} \uline{radiciel}, prenons $\deleted{\unsure{A'}}M'=G_{m\,k'}$, on a alors une suite exacte
% L3-4
\diagram{\begin{tikzcd}[ampersand replacement=\&,column sep=large]
\unsure{e|0}\arrow[r,no head]\&G_m\arrow[r]\&(G_m)_1\arrow[r]\arrow[d,equal]\&U\arrow[r]\&e\\
\&\&\underline{\operatorname{Hom}}_S(S',G_m)\&\&
\end{tikzcd}}
% L5-9
où $U$ est unipotent comm. \unsure{lisse|simple} sur $k$ [en car. 2, on peut trouver $U$ de dim 1, \deleted{\unsure{p.ex.}} si $[k':k]=2$]. La suite exacte de Coh. galoisienne sur $k$ donne alors, compte tenu de $H^{\unsure{i|1}}(k,(G_m)_1)\simeq H^i(k',G_m)$, l'isom
\[
H^1(k,U)\simeq\operatorname{Ker}(H^2(k,G_m)\unsure{{}-{ }|\longrightarrow}H^2(k',G_m))
\]

% L10-13
Prenons p.ex. $k=k_0[[t]]$, $k'=k_0[[t^{1/p}]]$, où $k_0$ est un corps fini de car. $p$, on trouve par le \unsure{corps de classes local|corps de classes locale}
\[
H^1(k,U)\simeq\mathbb{Z}/p\mathbb{Z}\qquad\ne0
\]
% slanted
\margin{N.B. On voit facilement que \unsure{ici on a} $pU=0$ (vrai \unsure{chaque fois} \unsure{pour|que} $k'^p\subset k$)}
% L14-15
[bien que, pour $p=2$, $U$ soit une forme de $G_a$ ; bien sur, cette forme n'est pas isom. à $G_a$].

% L16-19
Notons aussi que \uline{l'ext. de} $U$ \uline{par} $G_m$ [qui \unsure{splitte} sur $\bar k$] \uline{n'est pas triviale}, car autrement on aurait
\[
\underset{\substack{\shortparallel\\0}}{H^1(k,(G_m)_1)}\simeq\underset{\substack{\shortparallel\\0}}{H^1(k,G_m)}\times H^1(k,U)
\]
% L20-21
d'où $H^1(k,U)=0$. Donc ``\uline{le \unsure{r.|v.} \inserted{\uline{unipotent}} de} $(G_m)_1$ \uline{n'est pas défini sur} $k$''
